\documentstyle[%


righttag%


,11pt%


,amssymb%


,russian%


,izvru%


]{amsart}





\newcommand{\tl}[3]{\medskip\noindent{\bf#1}\ #2 \dotfill #3 \par} 


\pagestyle{empty}


\begin{document}


%\tableofcontents


\par


\vskip 2cm


\par


\bigskip


\bigskip


\bigskip


\par


\centerline{\Large СОДЕРЖАНИЕ}


\bigskip


\bigskip





\tl{Антонова О.В., Бронникова С.А.,


Давыдова В.В.,  Кабрера М.\,Ордоньез.}


{О слабом\newline законе больших чисел в банаховых


пространствах мартингального типа $p$ 


при общем условии\newline типа Чезаро}{3}


\tl{Акишев Г.А., Махашев С.Т.} 


{Об абсолютной сходимости рядов Фурье 


по обобщенной\newline системе Хаара}{8}


\tl{Арутюнянц Г.В.}


{Об ограниченности некоторых


максимальных и мультипликаторных\newline операторов}{17}


\tl{Дубровский В.В., Глекова В.В.} 


{Теорема о единственности решения обратных задач\newline


спектрального анализа}{21}


\tl{Ермолаева Л.Б.}


{Об одной квадратурной формуле}{25}


\tl{Захаров В.Н.} 


{Принцип абсолютного экстремума для одного 


класса дифференциальных\newline уравнений третьего 


порядка в трехмерном пространстве и его 


применение}{29}


\tl{Кириллов С.А., Кузнецов М.И., Чебочко Н.Г.} 


{О деформациях алгебры Ли типа $G_{2}$\newline характеристики три}{33}


\tl{Плещинский Н.Б.}


{К абстрактной теории приближенных методов 


решения линейных\newline операторных уравнений}{39}


\tl{Серебряков В.П.} 


{Об индексе дефекта матричных дифференциальных операторов 


второго\newline порядка с быстро осциллирующими коэффициентами}{48}


\tl{Федоров В.Е.} 


{Вырожденные сильно непрерывные


группы операторов}{54}


\tl{Ченцов А.Г.} 


{К вопросу об итерационной реализации неупреждающих многозначных\newline


отображений}{66}





\bigskip


{\centerline {КРАТКИЕ СООБЩЕНИЯ}}


\bigskip





\tl{Ахмадишина Ф.К.}


{Об асимптотической оптимальности 


проекционных методов}{77}


\tl{Молчанов А.В.} 


{Полугруппы эндоморфизмов слабых 


$p$-гиперграфов}{80}








\vfill


\noindent\raisebox{2pt}{\copyright} Казанский госудаpственный унивеpситет


\end{document}


